\documentclass[11pt, a4paper]{article}
\usepackage{geometry}
\usepackage[parfill]{parskip}
\usepackage{graphicx}
\usepackage{amsmath}
\usepackage{enumerate}

\usepackage{charter}

\title{Huiswerk 1 \\ IB0602 Lineaire algebra en stochastiek}
\author{Harman Singh}
\date{Month Day, Year}

\begin{document}
\maketitle

\pagebreak
\section*{Opgave 1}
\ldots



\pagebreak
\section*{Opgave 2}
\ldots




\pagebreak
\section*{Opgave 3}
\subsection*{oefening 1. $AB$}

Een $3\times2$ matrix vermenigvuldigd met een $2\times4$ matrix geeft een $3\times4$ matrix.

\[
\begin{pmatrix}
 (3)(4)+(1)(5) & (3)(1)+(1)(-2) & (3)(1)+(1)(1) & (3)(0)+(1)(3) \\
 (0)(4)+(-1)(5) & (0)(1)+(-1)(-2) & (0)(1)+(-1)(1) & (0)(0)+(-1)(3) \\
 (-2)(4)+(0)(5) & (-2)(1)+(0)(-2) & (-2)(1)+(0)(1) & (-2)(0)+(0)(3) \\
\end{pmatrix}
\]

\[
= \begin{pmatrix}
 (12+5) & (3-2) & (3+1) & (0+3) \\
 (0-5) & (0+2) & (0-1) & (0-3) \\
 (-8+0) & (-2+0) & (-2+0) & (0+0) \\
\end{pmatrix}
\]

\[
= \begin{pmatrix}
 17 & 1 & 4 & 3 \\
 -5 & 2 & -1 & -3 \\
 -8 & -2 & -2 & 0 \\
\end{pmatrix}
\]


\subsection*{oefening 2. $BA$}

Een $2\times4$ matrix vermenigvuldigd met een $3\times2$ matrix is niet mogelijk

\subsection*{oefening 3. ${B}^t{A}^t$}

Een $4\times2$ matrix vermenigvuldigd met een $2\times3$ matrix geeft een $4\times3$ matrix.

\[
\begin{pmatrix}
	(4)(3)+(5)(1) 	& (4)(0)+(5)(-1) 	& (4)(-2)+(5)(0) \\
	(1)(3)+(-2)(1) 	& (1)(0)+(-2)(-1) 	& (1)(-2)+(-2)(0) \\
	(1)(3)+(1)(1) 	& (1)(0)+(1)(-1) 	& (1)(-2)+(1)(0) \\
	(0)(3)+(3)(1) 	& (0)(0)+(3)(-1) 	& (0)(-2)+(3)(0) \\
\end{pmatrix}
\]

\[
= \begin{pmatrix}
	(12+5) 	& (0-5) & (-8+0) \\
	(3-2) 	& (0+2) & (-2+0) \\
	(3+1) 	& (0-1) & (-2+0) \\
	(0+3) 	& (0-3) & (0+0) \\
\end{pmatrix}
\]

\[
= \begin{pmatrix}
	17 & -5 & -8 \\
	1 & 2 & -2 \\
	4 & -1 & -2 \\
	3 & -3 & 0 \\
\end{pmatrix}
\]


\subsection*{oefening 3. $A{A}^t - 2{C}^t C$}





\ldots
\end{document}